% Acronyms used throughout the thesis.
\newacronym{acc}{ACC}{Acelerometria}
\newacronym{ai}{AI}{Artificial Intelligence (Inteligência Artificial)}
\newacronym{ann}{ANN}{Artificial Neural Network (Rede Neuronal Artificial)}
\newacronym{bic}{BIC}{Bayesian Information Criterion (Critério de Informação Bayesiano)}
\newacronym{ble}{BLE}{Bluetooth Low Energy}
\newacronym{bim}{BIM}{Building Information Modeling (Modelação da Informação da Construção)}
\newacronym{br}{BR}{Breathing Rate (Frequência Respiratória)}
\newacronym{cbt}{CBT}{Core Body Temperature (Temperatura Corporal Central)}
\newacronym{cnn}{CNN}{Convolutional Neural Network (Rede Neuronal Convolucional)}
\newacronym{csv}{CSV}{Comma-Separated Values}
\newacronym{criartE}{CRIARTE}{Construção com Robótica Inteligente e Arquitetura Revolucionária de Tecnologias Emergentes}
\newacronym{dbscan}{DBSCAN}{Density-Based Spatial Clustering of Applications with Noise}
\newacronym{dl}{DL}{Deep Learning (Aprendizagem Profunda)}
\newacronym{ecg}{ECG}{Eletrocardiograma}
\newacronym{eda}{EDA}{Electrodermal Activity (Atividade Eletrodérmica)}
\newacronym{eeg}{EEG}{Eletroencefalograma}
\newacronym{emg}{EMG}{Eletromiografia}
\newacronym{eog}{EOG}{Eletrooculografia}
\newacronym{fair}{FAIR}{Findable, Accessible, Interoperable and Reusable}
\newacronym{gmm}{GMM}{Gaussian Mixture Model (Modelo de Mistura Gaussiana)}
\newacronym{gnn}{GNN}{Graph Neural Network (Rede Neuronal de Grafos)}
\newacronym{gsr}{GSR}{Galvanic Skin Response (Resposta Galvânica da Pele)}
\newacronym{gecad}{GECAD}{Grupo de Engenharia e Computação Inteligente para a Inovação e o Desenvolvimento}
\newacronym{hr}{HR}{Heart Rate (Frequência Cardíaca)}
\newacronym{hrv}{HRV}{Heart Rate Variability (Variabilidade da Frequência Cardíaca)}
\newacronym{ia}{IA}{Inteligência Artificial}
\newacronym{imu}{IMU}{Inertial Measurement Unit (Unidade de Medição Inercial)}
\newacronym{iot}{IoT}{Internet of Things (Internet das Coisas)}
\newacronym{isep}{ISEP}{Instituto Superior de Engenharia do Porto}
\newacronym{iqr}{IQR}{Interquartile Range (Intervalo Interquartil)}
\newacronym{knn}{KNN}{K-Nearest Neighbors (K Vizinhos Mais Próximos)}
\newacronym{loso}{LOSO}{Leave-One-Subject-Out}
\newacronym{lstm}{LSTM}{Long Short-Term Memory}
\newacronym{ml}{ML}{Machine Learning (Aprendizagem Automática)}
\newacronym{met}{MET}{Metabolic Equivalent of Task (Equivalente Metabólico de Tarefa)}
\newacronym{nn}{NN}{Neural Network (Rede Neuronal)}
\newacronym{pca}{PCA}{Principal Component Analysis (Análise de Componentes Principais)}
\newacronym{pgf}{PGF}{Portable Graphics Format}
\newacronym{pnn}{pNN50}{Percentagem de intervalos NN com diferença superior a 50 ms}
\newacronym{ppg}{PPG}{Photoplethysmography (Fotopletismografia)}
\newacronym{prisma}{PRISMA}{Preferred Reporting Items for Systematic Reviews and Meta-Analyses}
\newacronym{prv}{PRV}{Pulse Rate Variability (Variabilidade do Ritmo de Pulso)}
\newacronym{rgpd}{RGPD}{Regulamento Geral sobre a Proteção de Dados}
\newacronym{rnn}{RNN}{Recurrent Neural Network (Rede Neuronal Recorrente)}
\newacronym{rr}{RR}{R-R Interval (Intervalo entre batimentos R-R)}
\newacronym{sdnn}{SDNN}{Standard Deviation of NN intervals (Desvio Padrão dos Intervalos NN)}
\newacronym{semg}{sEMG}{Surface Electromyography (Eletromiografia de Superfície)}
\newacronym{svm}{SVM}{Support Vector Machine (Máquina de Vetores de Suporte)}
\newacronym{svc}{SVC}{Support Vector Classification}
\newacronym{svr}{SVR}{Support Vector Regression}
\newacronym{ts}{TS}{Time Series (Série Temporal)}
\newacronym{vr}{VR}{Virtual Reality (Realidade Virtual)}
